% revised: CC 8 October 2026
% What this holds: the Data Appendix, through the end of the body.
% Read by arlington-bsap.tex. Prose lives here; formatting lives in the wrapper and style/.
% Not a part co-equal with A to C: a plain bold heading, like the Executive
% Summary, Works Cited and Legal Authorities, and its own headings
% below are not listed in the contents either (Sally, 7 October 2026).
\clearpage
\section*{Data Appendix}
\addstampedcontentsline{part}{Data Appendix}
\sectionbreaksfalse

\section*{Residents}

\subsection{Population}
Population is the decennial census count for the territory the Board governed, which
excludes the City of Alexandria throughout. Alexandria annexed part of the county in 1915
and again in 1930 (Districts, below), so part of the movement between the 1910 and 1930
censuses is territory changing hands, not people arriving or leaving.

\subsection{Race and Ethnicity}
From 1980 the groups in Figure~\ref{fig:residents-race} are the Census Bureau's race
categories crossed with Hispanic origin: Hispanic or Latino of any race, and every other
group non-Hispanic. Before 1980 the Bureau did not ask Hispanic origin of every household
and Hispanic residents were counted as white, so part of the step at 1980 is the Bureau starting to ask the question,
not people arriving. The 1980 Asian and Pacific Islander band also includes the
county's 384 American Indian, Eskimo and Aleut residents, because the 1980 table does not
separate them among residents of Spanish origin; the band is about five percent too high
in that one year.
``Other or Multiracial'' is American Indian and Alaska Native, some other race, and two or
more races.

\section*{Districts}

\subsection{District Lines}
The 1907 map of Alexandria County published under the Board of Supervisors' authority
is the only map found that draws the lines between the three magisterial districts; the
1879 and 1900 county maps letter each district's name across its interior and draw no
edge.\autocite{noetzel1907} Figure~\ref{fig:residents-district-map} reads each line off
the scan at about 15 points and joins them with straight segments, so its lines are
smoother than the printed ones. The scan is placed by fixing the sheet's corners to the
corners of the ten-mile square, and a line drawn this way is good to a few hundred
meters. The figure takes the 1907 lines to hold from 1870 to 1931: the census tables of
1880 to 1930 footnote no change to the districts other than Alexandria's two
annexations, and nothing found shows where the lines ran before 1880.

\subsection{The County's Outline}
The county's outline is today's, not 1907's: the County's own boundary file, joined to
the part of today's Alexandria, from the Census Bureau's, that lies inside the ten-mile
square. The two files are separate surveys and disagree by up to about 200 meters along
the edge they share; a strip narrower than 150 meters is dropped and a gap between them
is filled. The Potomac shore is today's, including land filled since 1907.

\subsection{Alexandria's Annexations}
The City of Alexandria's history of its boundaries maps each change on today's city
outline, and the figure takes the city as it stood before 1915, and the land annexed in
1915, from those maps, fitted to the city's outline to within about 13
meters.\autocite[3]{alexandria2024} The city drawn that way covers about 680 acres
against the 713 that Rose gives for 1912.\autocite{rose1964} Rose says the 1915
order carried the city's north line west to Braddock Road and gave the city only part of
the 866 acres it sought from the county; the city's map carries the north line west and
closes the area on a diagonal down to the old District line, 352 acres of county land.
That west edge rests on the city's drawing alone. The two sources disagree near Jones Point, where the city's map shows a wedge
annexed in 1915 that Rose's description of the 1912 limits puts inside the city; the
figure follows Rose. The city's map also sets the square's south corner at Jones Point
about 220 meters northeast of where Arlington's straight southwest edge, carried ten
miles, would put it, and the figure takes the city's position.

The 1930 annexation, which the city's history bounds by Duke Street, Quaker Lane and Four
Mile Run, is drawn as what is left of today's city inside the square once the 1912 city
and the 1915 area are out, about 2,660 acres. It therefore also holds any later change to
the boundary with Arlington, such as the small adjustments of 1966 along Walter Reed Drive
and Four Mile Run.\autocite[4]{alexandria2024} About 20 acres of it lie on Arlington
district's side of the 1907 line.

\section*{Elections}

\subsection{Votes for President}
The county's vote for President is the Commonwealth's official return by
county wherever one is printed. For 1876 through 1916 that is the
\textit{Warrock-Richardson Almanack}, a Richmond annual that prints "The
Official Vote of Virginia" from the official returns; from 1924 it is the
Secretary of the Commonwealth's own report, and from 1932 the Department of
Elections' database.\autocite{warrock1868,warrock1892,warrock1900,warrock1911,vasecretary1924,vasecretary1928,vaelections}
Two elections have no state return: for 1872 the county electoral board's own
certified table, printed in the \textit{Alexandria Gazette}, is used, and for
1920 Frank O'Leary's compilation of the \textit{Gazette}, the only count
held.\autocite{alexandriagazette1872official,oleary2010} For 1892 and 1900 the
\textit{Gazette}'s own district tables, summed, replace the Almanack's county
figure, which they contradict; in 1900 the two name different winners. The
Almanack is a compilation of the official returns and has not been checked
against the Secretary's own canvass for any year. Where the Almanack prints
only the leading tickets (1904 to 1916), minor-party votes are missing. The
county's candidate history prints the same returns from 1920 and agrees with
the state within five percent in every year but 1980, where the state's
canvass is kept.\autocite{arlingtonelections2021}

\subsection{Votes for County Board}
Before 1932 the Board's vote is the \textit{Alexandria Gazette}'s returns by
magisterial district in the issues after each election, read off the page
images, and the figures show it only in the years all three districts' counts
survive; each year's issue is cited in the archive delivered to the County.
From 1932 it is the Department of Elections'
locality rows, with the county's candidate history read as a check. A year with
two seats on the ballot records each voter twice, once per contest, so votes
per seat stand in for ballots cast in those years.
Figure~\ref{fig:election-sources} shows which publisher supplies the returns
for each year.

\begin{figure}[t]
  \centering
  \refstepcounter{figure}\label{fig:election-sources}
  \figuretitle{Election Returns by Publisher}
  \medskip
  \includegraphics[width=\textwidth]{elections_by_source.pdf}
  \smallskip
  \figurenotes{Notes: The dashed line marks 1932, when the Board went from three seats to five. A bar runs from the first to the last year a publisher
    supplies returns, and a row of ticks marks the scattered years where
    returns come from individual issues or pages. The press row gathers the
    \textit{Alexandria Gazette}, the \textit{Evening Star}, the \textit{Sun}
    and \textit{Daily Sun}, ARLnow and the Arlington GOP's candidate page.
    The county's candidate history also prints the presidential returns, which
    serve as a check on the state's.}
    {Source: \textit{Warrock-Richardson Almanack}; Frank O'Leary's electoral
    history; Secretary of the Commonwealth's reports; Arlington County's
    candidate history; Virginia Department of Elections database; the press
    items cited in the roster.}
\end{figure}

\section*{Board Members}

\subsection{Sources}
\begin{figure}[t]
  \centering
  \refstepcounter{figure}\label{fig:sources}
  \figuretitle{Board Member Records}
  \medskip
  \includegraphics[width=\textwidth]{members_by_source.pdf}
  \smallskip
  \figurenotes{Notes: The dashed line marks 1932, when the Board went from three seats to five. Each bar covers the years a term's source cites that
    document, so two bars over the same years are two documents confirming
    each other. The lighter colour marks a document that records members
    elected and the darker one a document that records members served; the
    Alexandria Gazette row carries both, since its items are a mix of the
    two. The Gazette row marks items the roster cites individually rather
    than as a term's primary source; single citations to the Washington
    Times, the Daily Sun and ARLnow are not shown.}
    {Source: Arlington Historical Society's 1967 compiled list of county
    officials; Frank O'Leary's electoral history; Norman Novack's roster;
    Arlington County's roll of members; the County's candidate history; the
    Virginia Department of Elections database; the Alexandria Gazette.}
\end{figure}

Several independent documents name who served on the Board across the years
since 1870, and their coverage overlaps rather than running end to end:
Figure~\ref{fig:sources} lays out the years each names, and where two
overlap, each confirms the other. The paragraphs below say what each
document gives, what it leaves out, and how the roster settles a disagreement
between two.

For 1870 through 1931 the roster rests on the Arlington Historical Society's 1967 list of county
officials. Its Research and Records Committee compiled the list from the Board's
own minute books, and it names each district's member term by
term.\autocite{arlhist1967officials} It dates a term from the first minute that
shows a man sitting, but it enters some departures only as ``Resigned; date ?''
and never states anyone's race. Frank O'Leary's electoral history, assembled from
the Alexandria Gazette and other newspapers, covers 1870 to 1920 and shows who won
each election.\autocite{oleary2010} It names a departure only where its prose
happens to, gives a surname alone for some members and initials alone before 1912.
Where the two disagree about who sat, the roster follows the minute books: an election return
names the winner, and the minute books name the man who took the seat. O'Leary
gives the Jefferson seat of 1870 to Storm V. Boyd on the election alone, but the
minute books say Boyd failed to qualify and never sat, so the seat stands vacant
until James C. Roach's appointment that September.

For 1932 through 1994 the roster rests on Norman Novack's list of members, printed in the Arlington
Historical Magazine in 1994.\autocite{novack1994} It records service rather than
contests, and it gives some dates to the day and others to the year alone. The
County's own roll of members, one entry for every year since 1932, dates arrivals
and departures to the day, including the members seated by appointment without an
election, and it is the second witness to Novack.\autocite{arlingtonva2026members}
The roll has slips of its own, among them duplicated names and a misspelling, so
it is read against Novack, and where the two disagree a third document decides. The
roll dates Howard Massey's appointment to November 10, 1952; Novack and the Daily
Sun's report of the three members a judge appointed that September date it
earlier, and the roster follows them.\autocite{dailysun1952appointees}

From 1995 on the roster rests on the County's roll, with the Office of Voter Registration and
Elections' candidate history, which runs to 2021, alongside it for the
contests.\autocite{arlingtonelections2021} The history lists candidates and
results, not the members seated between elections, so the roll is the only
document that shows an appointment such as Tannia Talento's on July 15, 2023. The
state's election results continue the contests after 2021.

A member's traits come from three kinds of source: the US census, which
describes members who served through the 1960s; the press, in news reports and
obituaries; and published profiles. A superscript on a value in Table~\ref{tab:roster}
marks its source: the census (c), the press or a published profile (p), or an
assumption (a) where no source says otherwise; two superscripts together mean
two kinds of source agree, and a dash means no source gives the value. The
census closes for 72 years, so the 1950 census is the last open.

\subsection{Gender}
Gender comes from two sources. Census sheets give it for members who served
through the 1960s. For nearly every other member the source is the press: the
``she'' or ``Ms.'' a contemporary article uses, quoted and cited member by
member. Neither source identifies Robinson and Willson, who served before 1900,
so both are assumed to be men.

\subsection{Race and Ethnicity}
Census records identify the five Black members of the 1870s and 1880s, and published
sources identify the six members of color seated since 1988.

Walter G. Willson is the one member of color identified only by a newspaper: the Alexandria
Gazette calls him ``colored'' in its reports of the elections of 1889 and 1891 and of his death in
1892, and no census record identifies him. The Gazette prints his name Wilson; the minute books and
O'Leary spell it Willson.\autocite{alexandriagazette18890525p3,alexandriagazette18910529p3,gazette1892wilsondeath}

Figure~\ref{fig:board-race-gender} carries one assumption: every member across 156 years
who is not named below as a member of color is recorded as White, because no source
says otherwise. For the 96 years between 1892 and 1988 the only evidence is one
sentence, the Arlington Historical Society's statement that William T. Newman,
Jr. was the first African American elected to the Board since Reconstruction;
no roster by race corroborates it.

For the six members before 1962 whose race no census or direct source gives,
a search of the Alexandria Gazette and the Evening Star by name found no racial
label. A missing label is weak evidence that a member was
White and no confirmation of it. The Gazette printed ``colored'' after Black
names in some reports, among them the Arlington Radicals' ticket of 1872, the
Alexandria city election of 1874 and the supervisors' elections of 1889 and
1891, but not in its county returns of 1879, which list Travis B. Pinn and
Tibbett Allen, both Black on the 1880 census, with no label, on a page whose
police report labels two other men.\autocite{alexandriagazette18720520p3,%
alexandriagazette18740529p2,alexandriagazette18890525p3,%
alexandriagazette18910529p3,gazette1879robinson,census1880pinn,census1880allen}
The Evening Star of 1957 named a person's race where race was the story, but
no Black candidate ran for the Board between 1958 and 1964, so nothing shows
whether the paper would have named Herbert L.\ Brown, Jr.\ or Thomas W.\
Richards as White.\autocite{star1957seating}

Figure~\ref{fig:board-race-coverage}, on page~\pageref{fig:board-race-coverage},
shows where the assumption matters most: most of the members it covers were
seated after 1962, when no census sheet is open to read.

\begin{figure}
  \centering
  \refstepcounter{figure}\label{fig:board-race-coverage}
  \figuretitle{Data Availability: Board Member Race}
  \medskip
  \includegraphics[width=\textwidth]{members_race_coverage.pdf}
\end{figure}

\subsection{Age}
Birth years come from census records for members seated before 1960 and from obituaries
and published profiles after that. Census birth years are derived from the age recorded when the census was taken, and are therefore accurate to within a year. No source gives a birth year for \birthYearMissingMembers{} of
\membersTotal{} members, who are left out of the figure.

Before 1932 the Board has three seats, and four members in that stretch have no birth year
in any source, so the age band there covers only the members with a birth year.
Medians are of ages on 1 July, and the county's adult median is interpolated within
census age bands.

\subsection{Residence}
% waits on: residence-county-records
\textit{[Where Board members have lived awaits the County's own records.]}

\subsection{Seat-Years}
From 1932 Figure~\ref{fig:board-race-gender} counts seats by the year, and a seat
whose holder changed during the year is split between the two. In 2023 Katie
Cristol left her seat in July and Tannia Talento was appointed to it, so that
year's band carries part of a seat for each, and the edge where one color meets
another sits inside the year. That is why the edges are ragged where a member
resigned, died or was appointed. Before 1932 each of the three district seats
is its own band, drawn to the month from the roster, so a change of member falls
where it happened and the first Board begins in July 1870.

In 11 years a seat stood empty while a special election was called: 1870, 1873,
1879 and 1920 in the district years, and 1990, 1997, 1999, 2003, 2012, 2014 and
2020 since; those years show a gap.

Table~\ref{tab:roster} lists every person who has served, in the order they
first took a seat, with the gender, race and birth year used in the figures. The
sources for every cell, one by one, are in the archive delivered to the
County.

% Generated by code/analysis/members_roster.py from data/clean/members.csv.
% Never edit by hand: bash run.sh rewrites it.
\clearpage\begingroup
\refstepcounter{table}\label{tab:roster}\xdef\rostertable{\thetable}
{\raggedright\MakeUppercase{Table \rostertable{}A}\par}
{\raggedright\bfseries Board Members, First Seated 1870--1899\par}
\medskip
\footnotesize\setlength{\tabcolsep}{8pt}\setlength{\LTpre}{0pt}\setlength{\LTpost}{0pt}
\begin{longtable}{Q}
\toprule \textit{Name} & \textit{Served} & \textit{Seated By} & \textit{Born} & \textit{Gender} & \textit{Race and Ethnicity} \\ \midrule \endfirsthead
\toprule \textit{Name} & \textit{Served} & \textit{Seated By} & \textit{Born} & \textit{Gender} & \textit{Race and Ethnicity} \\ \midrule \endhead
\bottomrule \endfoot
\multicolumn{6}{@{}l}{\bfseries 1870s} \\*
H. Dwight Smith & 1870--1873 & election & --- & man\textsuperscript{p} & White\textsuperscript{a} \\*
Charles W. Payne & 1870--1879 & election & 1822\textsuperscript{c} & man\textsuperscript{c} & White\textsuperscript{c} \\
James C. Roach & 1870--1871 & appointment & 1838\textsuperscript{c} & man\textsuperscript{c} & White\textsuperscript{c} \\
William A. Rowe & 1871--1883 & election & 1834\textsuperscript{c} & man\textsuperscript{c} & Black\textsuperscript{c,\,p} \\
Edward Deeble & 1871--1872 & election & 1811\textsuperscript{c} & man\textsuperscript{c} & White\textsuperscript{c} \\
John B. Syphax & 1872 & election & 1839\textsuperscript{c} & man\textsuperscript{c} & Black\textsuperscript{c,\,p} \\
Henry W. Febrey & 1872--1873 & election & 1828\textsuperscript{c} & man\textsuperscript{c} & White\textsuperscript{c} \\
Lott W. Crocker & 1873 & appointment & --- & man\textsuperscript{p} & White\textsuperscript{a} \\
Francis G. Schutt & 1873--1881 & appointment & 1834\textsuperscript{c} & man\textsuperscript{c,\,p} & White\textsuperscript{c} \\
Samuel Titus & 1873--1874 & appointment & 1815\textsuperscript{c} & man\textsuperscript{c} & White\textsuperscript{c} \\
Gilbert Vandenberg & 1874--1877 & election & 1820\textsuperscript{c} & man\textsuperscript{c} & White\textsuperscript{c} \\
William H. Robinson & 1877--1879 & election & --- & man\textsuperscript{a} & White\textsuperscript{a} \\*
Travis B. Pinn & 1879--1881 & election & 1841\textsuperscript{c} & man\textsuperscript{c} & Black\textsuperscript{c,\,p} \\
\addlinespace[1.4ex]
\multicolumn{6}{@{}l}{\bfseries 1880s} \\*
Christopher Costello & 1881--1885 & election & 1832\textsuperscript{c} & man\textsuperscript{c} & White\textsuperscript{c} \\*
Francis M. Mills & 1881--1883 & appointment & 1811\textsuperscript{c} & man\textsuperscript{c} & White\textsuperscript{c} \\
Perkins W. Squier & 1883--1884 & election & 1851\textsuperscript{c} & man\textsuperscript{c} & White\textsuperscript{c} \\
John W. Pendleton & 1883--1884 & election & 1848\textsuperscript{c} & man\textsuperscript{c} & Black\textsuperscript{c,\,p} \\
Curtis B. Graham, Jr. & 1884--1885 & appointment & 1849\textsuperscript{c} & man\textsuperscript{c} & White\textsuperscript{c} \\
George W. Saulisbury & 1884--1885 & --- & 1849\textsuperscript{c} & man\textsuperscript{c} & White\textsuperscript{c} \\
George W. Veitch & 1885--1887 & election & 1836\textsuperscript{c} & man\textsuperscript{c} & White\textsuperscript{c} \\
Richard W. Johnston & 1885--1887 & election & 1857\textsuperscript{c} & man\textsuperscript{c} & White\textsuperscript{c} \\
John D. Payne & 1885--1887 & election & 1855\textsuperscript{c} & man\textsuperscript{c} & White\textsuperscript{c} \\
Horatio Ball & 1887--1889 & election & 1827\textsuperscript{c} & man\textsuperscript{c} & White\textsuperscript{c} \\
Tibbett Allen & 1887--1888 & election & 1840\textsuperscript{c} & man\textsuperscript{c} & Black\textsuperscript{c,\,p} \\
A. B. Grunwell & 1887--1897 & election & 1834\textsuperscript{c} & man\textsuperscript{c} & White\textsuperscript{c} \\
Frank Hume & 1888--1895 & appointment & 1845\textsuperscript{c} & man\textsuperscript{c} & White\textsuperscript{c,\,p} \\
Frederick S. Corbett & 1889--1912 & election & 1856\textsuperscript{c} & man\textsuperscript{c} & White\textsuperscript{c} \\*
Walter G. Willson & 1889--1892 & election & --- & man\textsuperscript{a} & Black\textsuperscript{p} \\
\addlinespace[1.4ex]
\multicolumn{6}{@{}l}{\bfseries 1890s} \\*
Millard F. Birch & 1891--1893 & election & 1853\textsuperscript{c} & man\textsuperscript{c} & White\textsuperscript{c} \\*
William N. Febrey & 1892--1893 & appointment & 1851\textsuperscript{c} & man\textsuperscript{c,\,p} & White\textsuperscript{c} \\
John W. Clark & 1893--1895 & election & 1843\textsuperscript{c} & man\textsuperscript{c} & White\textsuperscript{c} \\
R. Henry Phillips & 1893--1895 & election & 1865\textsuperscript{p} & man\textsuperscript{p} & White\textsuperscript{a} \\
William Duncan & 1895--1904 & election & 1854\textsuperscript{c} & man\textsuperscript{c} & White\textsuperscript{c} \\
A. D. Torreyson & 1897 & election & 1866\textsuperscript{c} & man\textsuperscript{c} & White\textsuperscript{c} \\
George N. Saegmuller & 1897--1901 & election & 1847\textsuperscript{c} & man\textsuperscript{c} & White\textsuperscript{c} \\*
D. N. Rust & 1899--1908 & election & 1850\textsuperscript{c} & man\textsuperscript{c} & White\textsuperscript{c} \\
\end{longtable}\par\smallskip
\noindent\figurenotes{Notes: Superscripts indicate the source for each record: the US census (c), press or published profile (p), or assumption (a). Two superscripts mean two kinds of source agree. A dash means no source gives the value.}{}\endgroup
\clearpage\begingroup
{\raggedright\MakeUppercase{Table \rostertable{}B}\par}
{\raggedright\bfseries Board Members, First Seated 1900--1949\par}
\medskip
\footnotesize\setlength{\tabcolsep}{8pt}\setlength{\LTpre}{0pt}\setlength{\LTpost}{0pt}
\begin{longtable}{Q}
\toprule \textit{Name} & \textit{Served} & \textit{Seated By} & \textit{Born} & \textit{Gender} & \textit{Race and Ethnicity} \\ \midrule \endfirsthead
\toprule \textit{Name} & \textit{Served} & \textit{Seated By} & \textit{Born} & \textit{Gender} & \textit{Race and Ethnicity} \\ \midrule \endhead
\bottomrule \endfoot
\multicolumn{6}{@{}l}{\bfseries 1900s} \\*
Rezin W. Darbey & 1901--1904 & election & 1850\textsuperscript{c} & man\textsuperscript{c} & White\textsuperscript{c} \\*
Christopher J. Costello & 1901--1904 & election & 1865\textsuperscript{c} & man\textsuperscript{c} & White\textsuperscript{c} \\
W.W. Douglas & 1904--1908 & election & 1861\textsuperscript{c} & man\textsuperscript{c} & White\textsuperscript{c} \\
W. N. Febrey & 1904--1912 & election & 1851\textsuperscript{c} & man\textsuperscript{c} & White\textsuperscript{c} \\*
Edward Duncan & 1908--1932 & election & 1867\textsuperscript{p} & man\textsuperscript{c} & White\textsuperscript{c} \\
\addlinespace[1.4ex]
\multicolumn{6}{@{}l}{\bfseries 1910s} \\*
W. C. Wibirt & 1912--1919 & election & 1855\textsuperscript{c} & man\textsuperscript{c} & White\textsuperscript{c} \\*
Robert L. Walker & 1912--1919 & election & 1876\textsuperscript{c} & man\textsuperscript{c} & White\textsuperscript{c} \\*
Clarence R. Ahalt & 1919 & appointment & 1889\textsuperscript{c} & man\textsuperscript{c,\,p} & White\textsuperscript{c} \\
\addlinespace[1.4ex]
\multicolumn{6}{@{}l}{\bfseries 1920s} \\*
Thomas J. DeLashmutt & 1920--1923 & election & 1866\textsuperscript{c} & man\textsuperscript{c} & White\textsuperscript{c} \\*
Frank Upman & 1920 & appointment & 1873\textsuperscript{c} & man\textsuperscript{c} & White\textsuperscript{c} \\
W. T. Weaver & 1920--1923 & appointment & 1867\textsuperscript{c} & man\textsuperscript{c} & White\textsuperscript{c} \\
W. J. Ingram & 1924--1928 & election & 1867\textsuperscript{c} & man\textsuperscript{c} & White\textsuperscript{c} \\
E.C. Turnburke & 1924--1932 & election & 1864\textsuperscript{c} & man\textsuperscript{c} & White\textsuperscript{c} \\*
B. M. Hedrick & 1928--1932 & --- & 1897\textsuperscript{c} & man\textsuperscript{c} & White\textsuperscript{c} \\
\addlinespace[1.4ex]
\multicolumn{6}{@{}l}{\bfseries 1930s} \\*
Elizabeth B. Magruder & 1932--1947 & election & 1878\textsuperscript{c} & woman\textsuperscript{c,\,p} & White\textsuperscript{c} \\*
Fred A. Gosnell, Sr. & 1932--1933 & election & 1892\textsuperscript{c} & man\textsuperscript{c} & White\textsuperscript{c} \\
Harry A. Fellows & 1932--1935 & election & 1866\textsuperscript{c} & man\textsuperscript{c} & White\textsuperscript{c} \\
John C. Gall & 1932--1933 & election & 1901\textsuperscript{c} & man\textsuperscript{c} & White\textsuperscript{c} \\
Lyman M. Kelley & 1932--1934 & election & 1902\textsuperscript{c} & man\textsuperscript{c} & White\textsuperscript{c} \\
B. M. Smith & 1933--1935 & appointment & 1885\textsuperscript{c} & man\textsuperscript{c} & White\textsuperscript{c} \\
Christopher B. Garnett & 1933--1935 & appointment & 1876\textsuperscript{c} & man\textsuperscript{c} & White\textsuperscript{c} \\
W. A. E. McShea & 1934--1939 & appointment & 1865\textsuperscript{c} & man\textsuperscript{c} & White\textsuperscript{c} \\
F. Freeland Chew & 1936--1951 & election & 1898\textsuperscript{c} & man\textsuperscript{c} & White\textsuperscript{c} \\
George M. Yeatman & 1936--1939 & election & 1886\textsuperscript{c} & man\textsuperscript{c} & White\textsuperscript{c} \\*
W. P. Ames & 1936--1939 & election & 1894\textsuperscript{c} & man\textsuperscript{c,\,p} & White\textsuperscript{c} \\
\addlinespace[1.4ex]
\multicolumn{6}{@{}l}{\bfseries 1940s} \\*
Basil M. DeLashmutt & 1940--1949 & election & 1903\textsuperscript{c} & man\textsuperscript{c} & White\textsuperscript{c} \\*
Edmund D. Campbell & 1940--1946 & election & 1899\textsuperscript{c} & man\textsuperscript{c} & White\textsuperscript{c} \\
Leo C. Lloyd & 1940--1947 & election & 1890\textsuperscript{c} & man\textsuperscript{c} & White\textsuperscript{c} \\
Daniel A. Dugan & 1947--1952 & election & 1917\textsuperscript{c} & man\textsuperscript{c} & White\textsuperscript{c} \\
Harry W. Cuppett & 1947 & appointment & 1892\textsuperscript{c} & man\textsuperscript{c} & White\textsuperscript{c} \\
Alfred E. Frisbie & 1947--1952 & appointment & 1912\textsuperscript{c} & man\textsuperscript{c} & White\textsuperscript{c} \\*
Florence Cannon & 1948--1951 & election & 1889\textsuperscript{c} & woman\textsuperscript{c,\,p} & White\textsuperscript{c} \\
\end{longtable}\par\smallskip
\noindent\figurenotes{Notes: Superscripts indicate the source for each record: the US census (c), press or published profile (p), or assumption (a). Two superscripts mean two kinds of source agree. A dash means no source gives the value.}{}\endgroup
\clearpage\begingroup
{\raggedright\MakeUppercase{Table \rostertable{}C}\par}
{\raggedright\bfseries Board Members, First Seated 1950--1999\par}
\medskip
\footnotesize\setlength{\tabcolsep}{8pt}\setlength{\LTpre}{0pt}\setlength{\LTpost}{0pt}
\begin{longtable}{Q}
\toprule \textit{Name} & \textit{Served} & \textit{Seated By} & \textit{Born} & \textit{Gender} & \textit{Race and Ethnicity} \\ \midrule \endfirsthead
\toprule \textit{Name} & \textit{Served} & \textit{Seated By} & \textit{Born} & \textit{Gender} & \textit{Race and Ethnicity} \\ \midrule \endhead
\bottomrule \endfoot
\multicolumn{6}{@{}l}{\bfseries 1950s} \\*
Robert W. Cox & 1950--1952 & election & 1917\textsuperscript{c} & man\textsuperscript{c} & White\textsuperscript{c} \\*
Alan L. Dean & 1952 & election & 1919\textsuperscript{c} & man\textsuperscript{c} & White\textsuperscript{c} \\
Robert A. Peck & 1952--1955 & election & 1914\textsuperscript{c} & man\textsuperscript{c} & White\textsuperscript{c} \\
Howard R. Massey & 1952 & appointment & 1910\textsuperscript{c} & man\textsuperscript{c} & White\textsuperscript{c} \\
John A. Tillema & 1952 & appointment & 1887\textsuperscript{c} & man\textsuperscript{c} & White\textsuperscript{c} \\
M. Rex Byrne & 1952 & appointment & 1906\textsuperscript{c} & man\textsuperscript{c} & White\textsuperscript{c} \\
Alvin F. Kimel & 1952--1955 & special election & 1903\textsuperscript{c} & man\textsuperscript{c} & White\textsuperscript{c} \\
Leone B. Buchholz & 1952--1958 & special election & 1899\textsuperscript{c} & woman\textsuperscript{c,\,p} & White\textsuperscript{c} \\
Robert H. Detwiler & 1952--1953 & special election & 1910\textsuperscript{c} & man\textsuperscript{c} & White\textsuperscript{c} \\
George M. Rowzee, Jr. & 1953--1956 & election & 1913\textsuperscript{c} & man\textsuperscript{c} & White\textsuperscript{c} \\
Wesley W. Cooper & 1954--1957 & election & 1913\textsuperscript{c} & man\textsuperscript{c} & White\textsuperscript{c} \\
David L. Krupsaw & 1956--1960 & election & 1912\textsuperscript{c} & man\textsuperscript{c} & White\textsuperscript{c} \\
Ralph Kaul & 1956--1963 & election & 1914\textsuperscript{c} & man\textsuperscript{c} & White\textsuperscript{c} \\
Lucas H. Blevins & 1957--1960 & election & 1913\textsuperscript{c} & man\textsuperscript{c} & White\textsuperscript{c} \\
Herbert L. Brown, Jr. & 1958--1961 & election & 1913\textsuperscript{p} & man\textsuperscript{p} & White\textsuperscript{a} \\*
Leo Urbanske, Jr. & 1959--1966 & election & 1923\textsuperscript{c} & man\textsuperscript{c} & White\textsuperscript{c} \\
\addlinespace[1.4ex]
\multicolumn{6}{@{}l}{\bfseries 1960s} \\*
Ernest D. Wilt & 1960--1963 & appointment & 1911\textsuperscript{c} & man\textsuperscript{c} & White\textsuperscript{c} \\*
Thomas W. Richards & 1961--1975 & election & 1926\textsuperscript{p} & man\textsuperscript{p} & White\textsuperscript{a} \\
Roye L. Lowry & 1962--1965 & election & 1917\textsuperscript{p} & man\textsuperscript{p} & White\textsuperscript{a} \\
Harold J. Casto & 1964--1967 & election & 1924\textsuperscript{p} & man\textsuperscript{p} & White\textsuperscript{a} \\
Joseph L. Fisher & 1964--1974 & election & 1914\textsuperscript{p} & man\textsuperscript{p} & White\textsuperscript{a} \\
Dr. Kenneth M. Haggerty & 1966--1973 & election & 1924\textsuperscript{c} & man\textsuperscript{c} & White\textsuperscript{a} \\
Ned R. Thomas & 1967--1970 & election & 1929\textsuperscript{p} & man\textsuperscript{p} & White\textsuperscript{a} \\
Jay E. Ricks & 1968--1971 & election & 1932\textsuperscript{p} & man\textsuperscript{p} & White\textsuperscript{a} \\*
A. Leslie Phillips & 1969--1972 & election & 1921\textsuperscript{p} & man\textsuperscript{p} & White\textsuperscript{a} \\
\addlinespace[1.4ex]
\multicolumn{6}{@{}l}{\bfseries 1970s} \\*
Joseph S. Wholey & 1971--1978 & election & 1935\textsuperscript{p} & man\textsuperscript{p} & White\textsuperscript{a} \\*
Everard Munsey & 1972--1975 & election & 1933\textsuperscript{p} & man\textsuperscript{p} & White\textsuperscript{a} \\
John Purdy & 1973--1980 & election & 1931\textsuperscript{p} & man\textsuperscript{p} & White\textsuperscript{a} \\
Ellen Bozman & 1974--1997 & election & 1925\textsuperscript{p} & woman\textsuperscript{p} & White\textsuperscript{a} \\
Dorothy Grotos & 1976--1983 & election & 1931\textsuperscript{p} & woman\textsuperscript{p} & White\textsuperscript{a} \\
Walter Frankland & 1976--1983 & election & 1925\textsuperscript{p} & man\textsuperscript{p} & White\textsuperscript{a} \\*
Stephen Detwiler & 1979--1982 & election & 1943\textsuperscript{p} & man\textsuperscript{p} & White\textsuperscript{a} \\
\addlinespace[1.4ex]
\multicolumn{6}{@{}l}{\bfseries 1980s} \\*
John G. Milliken & 1981--1990 & election & 1946\textsuperscript{p} & man\textsuperscript{p} & White\textsuperscript{a} \\*
Mary Margaret Whipple & 1983--1995 & election & 1941\textsuperscript{p} & woman\textsuperscript{p} & White\textsuperscript{a} \\
Albert C. Eisenberg & 1984--1999 & election & 1946\textsuperscript{p} & man\textsuperscript{p} & White\textsuperscript{a} \\
Michael E. Brunner & 1984--1987 & election & 1943\textsuperscript{p} & man\textsuperscript{p} & White\textsuperscript{a} \\*
William T. Newman, Jr. & 1988--1993 & election & 1950\textsuperscript{p} & man\textsuperscript{p} & Black\textsuperscript{p} \\
\addlinespace[1.4ex]
\multicolumn{6}{@{}l}{\bfseries 1990s} \\*
James B. Hunter III & 1990--1997 & special election & 1940\textsuperscript{p} & man\textsuperscript{p} & White\textsuperscript{a} \\*
Benjamin H. Winslow, Jr. & 1993--1995 & special election & 1932\textsuperscript{p} & man\textsuperscript{p} & White\textsuperscript{a} \\
Christopher E. Zimmerman & 1996--2014 & special election & 1959\textsuperscript{p} & man\textsuperscript{p} & White\textsuperscript{a} \\
Paul F. Ferguson & 1996--2007 & election & 1965\textsuperscript{p} & man\textsuperscript{p} & White\textsuperscript{a} \\
Barbara A. Favola & 1997--2011 & special election & 1955\textsuperscript{p} & woman\textsuperscript{p} & White\textsuperscript{a} \\
G.N. “Jay” Fisette, Jr. & 1998--2017 & election & 1956\textsuperscript{p} & man\textsuperscript{p} & White\textsuperscript{a} \\*
Michael D. Lane & 1999 & special election & 1952\textsuperscript{p} & man\textsuperscript{p} & White\textsuperscript{a} \\
\end{longtable}\par\smallskip
\noindent\figurenotes{Notes: Superscripts indicate the source for each record: the US census (c), press or published profile (p), or assumption (a). Two superscripts mean two kinds of source agree. A dash means no source gives the value.}{}\endgroup
\clearpage\begingroup
{\raggedright\MakeUppercase{Table \rostertable{}D}\par}
{\raggedright\bfseries Board Members, First Seated 2000--2026\par}
\medskip
\footnotesize\setlength{\tabcolsep}{8pt}\setlength{\LTpre}{0pt}\setlength{\LTpost}{0pt}
\begin{longtable}{Q}
\toprule \textit{Name} & \textit{Served} & \textit{Seated By} & \textit{Born} & \textit{Gender} & \textit{Race and Ethnicity} \\ \midrule \endfirsthead
\toprule \textit{Name} & \textit{Served} & \textit{Seated By} & \textit{Born} & \textit{Gender} & \textit{Race and Ethnicity} \\ \midrule \endhead
\bottomrule \endfoot
\multicolumn{6}{@{}l}{\bfseries 2000s} \\*
Charles P. Monroe & 2000--2003 & election & 1957\textsuperscript{p} & man\textsuperscript{p} & Black\textsuperscript{p} \\*
J. Walter Tejada & 2003--2015 & special election & 1958\textsuperscript{p} & man\textsuperscript{p} & Hispanic\textsuperscript{p} \\*
Mary H. Hynes & 2008--2015 & election & 1956\textsuperscript{p} & woman\textsuperscript{p} & White\textsuperscript{a} \\
\addlinespace[1.4ex]
\multicolumn{6}{@{}l}{\bfseries 2010s} \\*
Libby T. Garvey & 2012--2024 & special election & 1951\textsuperscript{p} & woman\textsuperscript{p} & White\textsuperscript{a} \\*
John E. Vihstadt & 2014--2018 & special election & 1953\textsuperscript{p} & man\textsuperscript{p} & White\textsuperscript{a} \\
Christian E. Dorsey & 2016--2023 & election & 1972\textsuperscript{p} & man\textsuperscript{p} & Black\textsuperscript{p} \\
Kate A. "Katie" Cristol & 2016--2023 & election & 1985\textsuperscript{p} & woman\textsuperscript{p} & White\textsuperscript{a} \\
Erik Gutshall & 2018--2020 & election & 1970\textsuperscript{p} & man\textsuperscript{p} & White\textsuperscript{a} \\*
Matthew D. "Matt" de Ferranti & 2019-- & election & 1974\textsuperscript{p} & man\textsuperscript{p} & White\textsuperscript{a} \\
\addlinespace[1.4ex]
\multicolumn{6}{@{}l}{\bfseries 2020s} \\*
Takis P. Karantonis & 2020-- & special election & 1965\textsuperscript{p} & man\textsuperscript{p} & White\textsuperscript{a} \\*
Tannia Talento & 2023 & appointment & --- & woman\textsuperscript{p} & Hispanic\textsuperscript{p} \\
Maureen E. Coffey & 2024-- & election & 1995\textsuperscript{p} & woman\textsuperscript{p} & White\textsuperscript{a} \\
Susan R. Cunningham & 2024-- & election & --- & woman\textsuperscript{p} & White\textsuperscript{a} \\*
Julius D. "JD" Spain, Sr. & 2025-- & election & 1972\textsuperscript{p} & man\textsuperscript{p} & Black\textsuperscript{p} \\
\end{longtable}\par\smallskip
\noindent\figurenotes{Notes: Superscripts indicate the source for each record: the US census (c), press or published profile (p), or assumption (a). Two superscripts mean two kinds of source agree. A dash means no source gives the value.}{}\endgroup


\section*{Other Localities}

\subsection{Virginia}
Panel (a) of Figure~\ref{fig:localities-per-member} takes every Virginia city and
county of 100,000 or more residents in 2020. A city council's members are the
Richmond Charter Review Commission's count.\autocite[120--21]{richmond2023} A Mayor
elected at large counts as a member, because the commission's notes say most vote on
council; Richmond's is the one the notes place outside it, and is not counted. A
Mayor chosen from among the council is already one of its members. A county board's
members are read from the county's own page, and a chair elected at large counts as
one of them. Residents are the 2020 census count,\autocite{censusapi} and land area,
behind the density in Board Seats, is the Census Bureau's 2020
Gazetteer.\autocite{censusgazetteer2020}

\subsection{The Southeast}
Panel (b) takes every incorporated place, and every county or parish, of 150,000 to
300,000 residents in the 2020 census in the nine states the commission's Appendix E
compares (Alabama, Arkansas, Georgia, Louisiana, Mississippi, North Carolina, South
Carolina, Tennessee and Virginia) and in Maryland, which has no incorporated place
that size and enters by county. The states are the commission's; the floor and the
counties are this report's, because the commission's 180,000 floor and its list of
cities alone answer Richmond's question, and Arlington governs as a county. A county
is left out where the Census Bureau's county is not the body that governs it:
Virginia's independent cities are counted as cities, and three Georgia counties are
the consolidated governments of Augusta, Columbus and Macon-Bibb County, counted once,
as those cities.

The 19 cities Appendix E tabulates take its count of council
members,\autocite[122]{richmond2023} a Mayor who sits and votes counted as one seat;
this is why its Norfolk has eight members where Appendix D's count of seven leaves
out the Mayor elected at large. Every other body's members are read from its own
charter or official page, one source to a body, and the same rule holds throughout:
every member who sits and votes counts, a chair or Mayor elected at large included,
and a chief executive who does not sit, or votes only to break a tie, does not.
Where a page lists members without saying how the body is made up, the count is the
roster's. Residents are the 2020 census count for the place or county, not the
source's.

Three counts could change. Appendix E states no date for its table. Clarksville,
Tennessee, is counted as twelve ward seats, though a 2009 opinion of Tennessee's
Attorney General reads its charter as giving the Mayor a vote, which would make
thirteen. St.~Tammany Parish's council may change after a special election on
November 3, 2026.

